\documentclass[12pt,letterpaper,titlepage]{article}

\usepackage{fontspec}
\defaultfontfeatures{Mapping=tex-text}
\usepackage{xunicode}
\usepackage{xltxtra}
\usepackage{amsmath}
\usepackage{pdfpages}
\usepackage{mathtools}
\usepackage{amsfonts}
\usepackage{amssymb}
\setcounter{secnumdepth}{0}
\usepackage{nameref}
\usepackage{enumitem}
\usepackage{environ}
\usepackage{mathtools}
\DeclarePairedDelimiter{\ceil}{\lceil}{\rceil}

\showboxdepth=\maxdimen
\showboxbreadth=\maxdimen

\usepackage[tocflat]{tocstyle}
\usetocstyle{allwithdot}
\usepackage[bottom]{footmisc}

\usepackage{paracol}
\usepackage{wrapfig}
\globalcounter{table}
\globalcounter{figure}
\usepackage{graphicx}
\usepackage[left=1in,right=1in,top=1in,bottom=1in]{geometry}
\graphicspath{{img/}}

\author{Jacob Abel}
\title{	Homework 10
	\\\large MATH2534 CRN:15708
}

\setlength{\parskip}{0.5em}

\begin{document}
\maketitle
\begin{raggedright}

\begin{enumerate}

\item Let $X = \{1, 2, 3, 4\}$ and $Y = \{a, b, c, d, e\}$. Define $f : X \mapsto Y$ as follows:
\begin{equation*}
f(1) = c,\;f(2) = a,\;f(3) = e,\;\text{and}\;f(4) = d
\end{equation*}
\begin{enumerate}[label=(\alph*)]
\item Draw an arrow diagram and determine if $f$ is a function.
\begin{center}
\includegraphics[width=.30\textwidth, height=\textheight, keepaspectratio=true]{hw10q1a}
\end{center}

\item Find the domain, codomain and range of $f$.

Domain: $X$, Codomain: $Y$, Range: $\{a, c, d, e\}$
\item Let $A = \{b\}$, $B = \{a, d\}$, and $C = \{b, e\}$. The inverse image of a set $S$ is the set of all preimages of each element of $S$. Find the inverse images of $A$, $B$, and $C$.
$f^{-1}(A)=\emptyset$, $f^{-1}(B)=\{2, 4\}$, $f^{-1}(C)=\{3\}$
\end{enumerate}

\item Let $f : \mathbb{Q} \mapsto \mathbb{Q}$ be given by $f(\frac{m}{n}) = m + n$ for all integers $m$ and $n$ with $n \neq 0$. Is $f$ well-defined? Why or why not?

For $n = 0$, $f$ is undefined. $n = 0$ is outside of the defined boundaries of $f$. Therefore $f$ is well defined for all integers $m$ and $n$ with $n \neq 0$.

\item Let $A = \{1, 2, 3, 4, 5\}$ and define $S : \mathcal{P}(A) \mapsto Z$ as follows: For all sets $X$ in $\mathcal{P}(A)$,

\begin{equation*}
S(X) = \begin{cases} 
      1 & \text{if}\, A − X\, \text{has an even number of elements} \\
      0 & \text{if}\, A − X\, \text{has an odd number of elements}
   \end{cases}
\end{equation*}

\begin{enumerate}[label=(\alph*)]
\item Find $S(\emptyset)$, $S(\{1, 4\})$, and $S(\{1, 2, 3, 4, 5\})$.

$S(\emptyset) = 1$

$S(\{1, 4\}) = 1$

$S(\{1, 2, 3, 4, 5\}) = 0$
\item Is $S$ a function? Why or why not?

$S$ is a function as it always cleanly maps from any $X$ into some value $Z$.
\end{enumerate}

\clearpage

\item Let $A = \{a_1, a_2, a_3\}$ and $\{0, 1\}^3 = \{0, 1\} \times \{0, 1\} \times \{0, 1\}$ be the set of all ordered 3-tuples of 0’s and 1’s. Define a function $f : P(A) \mapsto \{0, 1\}^3$ by $f(X) = (s_1, s_2, s_3)$ for $X \in P(A)$ where

\begin{equation*}
s_i = \begin{cases} 
      1 & \text{if}\, a_i\in X \\
      0 & \text{if}\, a_i\not\in X
   \end{cases}
\end{equation*}
Draw an arrow diagram for $f$.

\underline{Ran out of time. Did not complete.}

\item Determine whether the relation given is one-to-one, onto. If not, explain why not.

\begin{enumerate}[label=(\alph*)]
\item Let $f : R \mapsto Z$ be defined by $f(x) = \ceil[\big]{x}$.

It is onto as every rational number can be rounded up to some integer. 

It is not one-to-one as multiple rational numbers can round up to an integer.

\item Let $Y = \{0, 1, 2, 3, 4, 5\}$ where $f : Z^+\mapsto Y$ is defined by $f(x) = 5x \text{mod} 6$.

It is onto as every positive integer $x$ mod 6 will map onto some value in $Y$. This is not changed what that integer $x$ is multiplied by 5.

It is not one-to-one as every 1 out of 6 positive integers maps to each value in $Y$.

\item Let $M : Z^+ \times Z^+ \mapsto Z^+$ be given by $M(m, n) = m + n$.

It is onto as any addition between 2 positive integers will always result in another positive integer.

It is not one-to-one as a variety of different integers can be added together to get the same output integer. i.e. $1+3 = 4, 2+2 = 4$

\end{enumerate}

\item Determine whether the following mapping is one-to-one, onto, bijective? Find an inverse function if it exists. Justify your answer.

\begin{enumerate}[label=(\alph*)]
\item Let $f : R \mapsto R$ be defined by $f(x) = −2x + 6$.

The function is bijective as every input has 1 and only 1 output and vice versa. (1 Variable, Linear function)

The inverse,  $f^{-1} : R \mapsto R$ is defined as $f^{-1}(x) = 3 - \frac{y}{2}$

\item Let $f : R − \{3\} \mapsto R$ be defined by $f(x) = \frac{2x}{x − 3}$.

This function is bijective as every value in $R-\{3\}$ maps directly and individually onto some value in $R$ and vice versa.

The inverse, $f^{-1} : R − \{2\} \mapsto R$ is defined as $f^{-1}(x) = \frac{3y}{y-2}$.

\end{enumerate}

\item Prove the statement:

If $f : X \mapsto Y$ and $g : Y \mapsto Z$, are both one-to-one functions, then $g \circ f : X \mapsto Z$ is one-to-one.

\underline{Ran out of time. Did not complete.}
\end{enumerate}
\end{raggedright}
\end{document}
